\documentclass[10pt,a4paper]{amsart}
\usepackage{amsmath,amssymb,amsthm,graphicx}
\usepackage{verbatim}
\newenvironment{reference}{\comment}{\endcomment}
%\newenvironment{reference}{}{}
\newenvironment{slogan}{\comment}{\endcomment}
\newenvironment{history}{\comment}{\endcomment}

% For commutative diagrams we use Xy-pic
\usepackage[all]{xy}

% We use 2cell for 2-commutative diagrams.
\xyoption{2cell}
% Theorem environments.
%
\theoremstyle{plain}
\newtheorem{theorem}[subsection]{Theorem}
\newtheorem{proposition}[subsection]{Proposition}
\newtheorem{lemma}[subsection]{Lemma}

\theoremstyle{definition}
\newtheorem{definition}[subsection]{Definition}
\newtheorem{example}[subsection]{Example}
\newtheorem{exercise}[subsection]{Exercise}
\newtheorem{situation}[subsection]{Situation}

\theoremstyle{remark}
\newtheorem{remark}[subsection]{Remark}
\newtheorem{remarks}[subsection]{Remarks}

\numberwithin{equation}{subsection}

% Macros
%
\def\lim{\mathop{\mathrm{lim}}\nolimits}
\def\colim{\mathop{\mathrm{colim}}\nolimits}
\def\Spec{\mathop{\mathrm{Spec}}}
\def\Hom{\mathop{\mathrm{Hom}}\nolimits}
\def\Ext{\mathop{\mathrm{Ext}}\nolimits}
\def\SheafHom{\mathop{\mathcal{H}\!\mathit{om}}\nolimits}
\def\SheafExt{\mathop{\mathcal{E}\!\mathit{xt}}\nolimits}
\def\Sch{\mathit{Sch}}
\def\Mor{\mathop{\mathrm{Mor}}\nolimits}
\def\Ob{\mathop{\mathrm{Ob}}\nolimits}
\def\Sh{\mathop{\mathit{Sh}}\nolimits}
\def\NL{\mathop{N\!L}\nolimits}
\def\CH{\mathop{\mathrm{CH}}\nolimits}
\def\proetale{{pro\text{-}\acute{e}tale}}
\def\etale{{\acute{e}tale}}
\def\QCoh{\mathit{QCoh}}
\def\Ker{\mathop{\mathrm{Ker}}}
\def\Im{\mathop{\mathrm{Im}}}
\def\Coker{\mathop{\mathrm{Coker}}}
\def\Coim{\mathop{\mathrm{Coim}}}

% Boxtimes
%
\DeclareMathSymbol{\boxtimes}{\mathbin}{AMSa}{"02}

%
% Macros for moduli stacks/spaces
%
\def\QCohstack{\mathcal{QC}\!\mathit{oh}}
\def\Cohstack{\mathcal{C}\!\mathit{oh}}
\def\Spacesstack{\mathcal{S}\!\mathit{paces}}
\def\Quotfunctor{\mathrm{Quot}}
\def\Hilbfunctor{\mathrm{Hilb}}
\def\Curvesstack{\mathcal{C}\!\mathit{urves}}
\def\Polarizedstack{\mathcal{P}\!\mathit{olarized}}
\def\Complexesstack{\mathcal{C}\!\mathit{omplexes}}
% \Pic is the operator that assigns to X its picard group, usage \Pic(X)
% \Picardstack_{X/B} denotes the Picard stack of X over B
% \Picardfunctor_{X/B} denotes the Picard functor of X over B
\def\Pic{\mathop{\mathrm{Pic}}\nolimits}
\def\Picardstack{\mathcal{P}\!\mathit{ic}}
\def\Picardfunctor{\mathrm{Pic}}
\def\Deformationcategory{\mathcal{D}\!\mathit{ef}}
\usepackage{cleveref}
\crefname{originallemma}{original lemma}{original lemmas}
\Crefname{originallemma}{Original lemma}{Original lemmas}
\makeatletter
\newlabel{lemma-dimension-function-scheme}{{0AWF}{source}{}{sourcepointer-1}{}}
\newlabel{lemma-dimension-function-scheme@cref}{{[originallemma][0][]0AWF}{[1][0][]source}{}{}{}}
\makeatother
\crefname{originallemma}{original lemma}{original lemmas}
\Crefname{originallemma}{Original lemma}{Original lemmas}
\makeatletter
\newlabel{morphisms-lemma-dimension-function-propagates}{{02JW}{source}{}{sourcepointer-2}{}}
\newlabel{morphisms-lemma-dimension-function-propagates@cref}{{[originallemma][0][]02JW}{[1][0][]source}{}{}{}}
\makeatother
\crefname{originallemma}{original lemma}{original lemmas}
\Crefname{originallemma}{Original lemma}{Original lemmas}
\makeatletter
\newlabel{topology-lemma-dimension-function-unique}{{02IB}{source}{}{sourcepointer-3}{}}
\newlabel{topology-lemma-dimension-function-unique@cref}{{[originallemma][0][]02IB}{[1][0][]source}{}{}{}}
\makeatother
\crefname{originallemma}{original lemma}{original lemmas}
\Crefname{originallemma}{Original lemma}{Original lemmas}
\makeatletter
\newlabel{properties-lemma-locally-Noetherian-closed-point}{{02IL}{source}{}{sourcepointer-4}{}}
\newlabel{properties-lemma-locally-Noetherian-closed-point@cref}{{[originallemma][0][]02IL}{[1][0][]source}{}{}{}}
\makeatother
\crefname{originallemma}{original lemma}{original lemmas}
\Crefname{originallemma}{Original lemma}{Original lemmas}
\makeatletter
\newlabel{lemma-good-dualizing-normalized}{{0AWL}{source}{}{sourcepointer-5}{}}
\newlabel{lemma-good-dualizing-normalized@cref}{{[originallemma][0][]0AWL}{[1][0][]source}{}{}{}}
\makeatother
\title{The dimension function of the normalised dualizing complex differs from the base one by the transcendence degree}
\author{The Stacks Project Authors}
\date{stacks-project:duality.tex:local-snapshot-20261012; GNU FDL 1.2 or later}
\begin{document}
\maketitle
\noindent\textit{Editorial note.} All delivered mathematical text is the author's own, copied from the pinned Stacks Project chapter file duality.tex; no statement or proof step was rewritten, reordered or removed.\par\medskip
\noindent\textit{Editorial note.} Preamble provenance: the pinned chapter file opens with an input of the shared preamble and carries no preamble text of its own, so the project's single Stacks-wide preamble staged at sources/stacks-duality/preamble.tex is excerpted verbatim under the recipe fixed by LANE\_BRIEF 22.1, as three spans, lines 8-13 (verbatim package and the reference/slogan/history environments), lines 14-18 (Xy-pic setup) and lines 154-217 (the native theorem declarations and mathematical macros). Nothing was added, removed or redefined; the wrapper class and the amsmath/amssymb/amsthm/graphicx base packages are supplied by the assembler.\par\medskip
\noindent\textit{Editorial note.} The counted statement and proof contain no cross reference, no citation and no included file, so the delivered body is self contained.\par\medskip
\noindent\textit{Editorial note.} Canonical location of the counted proposition: Stacks tag 0AWM (sources/stacks-duality/pinned-original-source.tex line 5856).\par\medskip
\noindent\textit{Allowed standard prerequisites (editorial).}\begin{itemize}
\item The dimension function attached to a dualizing complex and universal catenarity. Same chapter, Lemma labelled lemma-dimension-function-scheme, cited at proof lines 5908-5910 and in the statement. Supplies the dimension functions and the catenarity of S.
\item Dimension functions propagate and are unique up to a locally constant difference. Morphisms, Lemma labelled morphisms-lemma-dimension-function-propagates, cited at proof lines 5904-5908, and Topology, Lemma labelled topology-lemma-dimension-function-unique, cited at proof lines 5910-5912. Reduce the claim to a single point.
\item Existence of a closed point and normalisation at closed points. Properties, Lemma labelled properties-lemma-locally-Noetherian-closed-point, cited at proof lines 5914-5916, and Same chapter, Lemma labelled lemma-good-dualizing-normalized, cited at proof lines 5916. Gives the point where the equality is checked.
\end{itemize}\medskip
\ifdefined\hypertarget\hypertarget{sourcepointer-1}{}\fi
\noindent\textit{Source reference.} Original lemma 0AWF: sources/stacks-duality/pinned-original-source.tex, source lines 567--567. Same-chapter pointer into the author environment lemma at source line 567; invoked at proof line 5881 of the counted proof. It is retained only as a navigation pointer; every load-bearing step is inside the delivered span.\par\medskip
\ifdefined\hypertarget\hypertarget{sourcepointer-2}{}\fi
\noindent\textit{Source reference.} Original lemma 02JW: sources/stacks-chow/tags.txt, source lines 3351--3351. Cross-chapter standard prerequisite located in the project tag list staged in the local snapshot; invoked at proof line 5879 of the counted proof. It is retained only as a navigation pointer; every load-bearing step is inside the delivered span.\par\medskip
\ifdefined\hypertarget\hypertarget{sourcepointer-3}{}\fi
\noindent\textit{Source reference.} Original lemma 02IB: sources/stacks-chow/tags.txt, source lines 3294--3294. Cross-chapter standard prerequisite located in the project tag list staged in the local snapshot; invoked at proof line 5882 of the counted proof. It is retained only as a navigation pointer; every load-bearing step is inside the delivered span.\par\medskip
\ifdefined\hypertarget\hypertarget{sourcepointer-4}{}\fi
\noindent\textit{Source reference.} Original lemma 02IL: sources/stacks-chow/tags.txt, source lines 3304--3304. Cross-chapter standard prerequisite located in the project tag list staged in the local snapshot; invoked at proof line 5885 of the counted proof. It is retained only as a navigation pointer; every load-bearing step is inside the delivered span.\par\medskip
\ifdefined\hypertarget\hypertarget{sourcepointer-5}{}\fi
\noindent\textit{Source reference.} Original lemma 0AWL: sources/stacks-duality/pinned-original-source.tex, source lines 5805--5805. Same-chapter pointer into the author environment lemma at source line 5805; invoked at proof line 5888 of the counted proof. It is retained only as a navigation pointer; every load-bearing step is inside the delivered span.\par\medskip
\begin{lemma}
\label{lemma-good-dualizing-dimension-function}
Let $S$ be a Noetherian scheme and let $\omega_S^\bullet$ be a
dualizing complex. Let $f : X \to S$ be of finite type
and let $\omega_X^\bullet$ be the dualizing complex
normalized relative to $\omega_S^\bullet$. For all $x \in X$ we have
$$
\delta_X(x) - \delta_S(f(x)) = \text{trdeg}_{\kappa(f(x))}(\kappa(x))
$$
where $\delta_S$, resp.\ $\delta_X$
is the dimension function of
$\omega_S^\bullet$, resp.\ $\omega_X^\bullet$, see
Lemma \ref{lemma-dimension-function-scheme}.
\end{lemma}

\begin{proof}
We may replace $X$ by an affine neighbourhood of $x$. Hence we may
and do assume there is a compactification $X \subset \overline{X}$
over $S$. Then we may replace $X$ by $\overline{X}$ and assume
that $X$ is proper over $S$. We may also assume $X$ is connected
by replacing $X$ by the connected component of $X$ containing $x$.
Next, recall that both $\delta_X$ and the function
$x \mapsto \delta_S(f(x)) + \text{trdeg}_{\kappa(f(x))}(\kappa(x))$
are dimension functions on $X$, see
Morphisms, Lemma \ref{morphisms-lemma-dimension-function-propagates}
(and the fact that $S$ is universally catenary by
Lemma \ref{lemma-dimension-function-scheme}).
By Topology, Lemma \ref{topology-lemma-dimension-function-unique}
we see that the difference is locally constant, hence constant as $X$ is
connected. Thus it suffices to prove equality in any point of $X$.
By Properties, Lemma \ref{properties-lemma-locally-Noetherian-closed-point}
the scheme $X$ has a closed point $x$. Since $X \to S$ is proper
the image $s$ of $x$ is closed in $S$. Thus we may apply
Lemma \ref{lemma-good-dualizing-normalized} to conclude.
\end{proof}

\end{document}
